% AUTO-GENERATED by the update_record tool. DO NOT EDIT.
% verify_paper regenerates this file from record.json and the run outputs
% and fails on any difference, so manual edits always surface.
\noindent\textbf{H1.} SciGym 論文の Table 1 のうち Gemini-2.5-Pro、Gemini-2.5-Flash、GPT-4.1、GPT-4.1-mini の 4 行は、公開コードの Controller をそのまま用い、現在 API で提供されている同名モデル（Gemini 2 種は GA 版、GPT-4.1 2 種は論文と同じ 2025-04-14 スナップショット）で SciGym-small 137 件を最大 20 反復で解かせ、提出モデルを airas-eval の scigym\_small パック（SciGym 公式 Evaluator）で採点したとき、各モデルの平均 Simulation Trajectory Error（STE）が論文値から 0.05 以内に収まる形で再現される。
\emph{Grounded on:} \texttt{\detokenize{s1.p1}}, \texttt{\detokenize{s1.p5}}, \texttt{\detokenize{s1.p10}}.
\begin{enumerate}
\item[\textbf{C1}] 現行の Gemini-2.5-Pro（google/gemini-2.5-pro）で SciGym-small 137 件を最大 20 反復で解かせたときの平均 STE は、論文の 0.3212 を 0.05 より大きくは上回らない。
  \emph{Rationale:} Table 1 で最良とされた行がそのまま再現されることは、論文の主結果が特定のプレビュー版や実行環境に依存していないことの直接の証拠であり、h1 の Gemini-2.5-Pro 行を担う。
  \emph{Cites:} \texttt{\detokenize{s1.p5}}, \texttt{\detokenize{s1.p10}}, \texttt{\detokenize{s1.p3}}, \texttt{\detokenize{s1.p2}}, \texttt{\detokenize{s1.p7}}; criterion: \texttt{\detokenize{s1.p5}}; \texttt{\detokenize{d1}}: \texttt{\detokenize{s1.p3}}, \texttt{\detokenize{s1.p4}}, \texttt{\detokenize{s2.p1}}, \texttt{\detokenize{s3.p1}}; \texttt{\detokenize{gemini-2.5-pro}}: \texttt{\detokenize{s1.p1}}.
  \emph{Criterion:} \texttt{\detokenize{gemini-2.5-pro}}.\texttt{\detokenize{trajectory_smape}} $-$ 0.3212 $\leq 0.05$.
  \emph{Prediction:} $[-0.05, 0.05]$ (論文 Table 1 の 1 回実行の平均値。前回の再現（auto-res2/scigym-table1-reproduction）で 1 件の STE の標準偏差は約 0.3、137 件平均の標準誤差は約 0.026 だった。モデル版の差を含めても ±0.05 に収まると予測する).
  \emph{Observed:} pending.
  \emph{Verdict:} pending.
\item[\textbf{C2}] 現行の Gemini-2.5-Flash（google/gemini-2.5-flash）で SciGym-small 137 件を最大 20 反復で解かせたときの平均 STE は、論文の 0.4181 を 0.05 より大きくは上回らない。
  \emph{Rationale:} Gemini 系の mini 変種の行が再現されることで、h1 の Gemini-2.5-Flash 行を担い、c1 と合わせて Gemini 系の pro / mini の関係が論文どおりかを読めるようにする。
  \emph{Cites:} \texttt{\detokenize{s1.p5}}, \texttt{\detokenize{s1.p10}}, \texttt{\detokenize{s1.p3}}, \texttt{\detokenize{s1.p2}}, \texttt{\detokenize{s1.p6}}; criterion: \texttt{\detokenize{s1.p5}}; \texttt{\detokenize{d2}}: \texttt{\detokenize{s1.p3}}, \texttt{\detokenize{s1.p4}}, \texttt{\detokenize{s2.p1}}, \texttt{\detokenize{s3.p1}}; \texttt{\detokenize{gemini-2.5-flash}}: \texttt{\detokenize{s1.p1}}.
  \emph{Criterion:} \texttt{\detokenize{gemini-2.5-flash}}.\texttt{\detokenize{trajectory_smape}} $-$ 0.4181 $\leq 0.05$.
  \emph{Prediction:} $[-0.05, 0.05]$ (論文 Table 1 の 1 回実行の平均値。前回の再現での差は +0.028 で、137 件平均の標準誤差は約 0.028 だった).
  \emph{Observed:} pending.
  \emph{Verdict:} pending.
\item[\textbf{C3}] GPT-4.1（openai/gpt-4.1、2025-04-14 スナップショット）で SciGym-small 137 件を最大 20 反復で解かせたときの平均 STE は、論文の 0.4611 を 0.05 より大きくは上回らない。
  \emph{Rationale:} 論文と同じ日付スナップショットが使える唯一の pro モデルであり、モデル版の差を含まない最も純粋な再現条件として h1 の GPT-4.1 行を担う。
  \emph{Cites:} \texttt{\detokenize{s1.p5}}, \texttt{\detokenize{s1.p10}}, \texttt{\detokenize{s1.p3}}, \texttt{\detokenize{s1.p2}}; criterion: \texttt{\detokenize{s1.p5}}; \texttt{\detokenize{d3}}: \texttt{\detokenize{s1.p3}}, \texttt{\detokenize{s1.p4}}, \texttt{\detokenize{s2.p1}}, \texttt{\detokenize{s3.p1}}; \texttt{\detokenize{gpt-4.1}}: \texttt{\detokenize{s1.p1}}.
  \emph{Criterion:} \texttt{\detokenize{gpt-4.1}}.\texttt{\detokenize{trajectory_smape}} $-$ 0.4611 $\leq 0.05$.
  \emph{Prediction:} $[-0.05, 0.05]$ (論文 Table 1 の 1 回実行の平均値。前回の再現での差は +0.047 で、137 件平均の標準誤差は約 0.023 だった).
  \emph{Observed:} pending.
  \emph{Verdict:} pending.
\item[\textbf{C4}] GPT-4.1-mini（openai/gpt-4.1-mini、2025-04-14 スナップショット）で SciGym-small 137 件を最大 20 反復で解かせたときの平均 STE は、論文の 0.6007 を 0.05 より大きくは上回らない。
  \emph{Rationale:} 同じスナップショットが使える mini モデルの行を担い、c3 と合わせて OpenAI 系の pro / mini の関係が論文どおりかを読めるようにする。
  \emph{Cites:} \texttt{\detokenize{s1.p5}}, \texttt{\detokenize{s1.p10}}, \texttt{\detokenize{s1.p3}}, \texttt{\detokenize{s1.p2}}, \texttt{\detokenize{s1.p6}}; criterion: \texttt{\detokenize{s1.p5}}; \texttt{\detokenize{d4}}: \texttt{\detokenize{s1.p3}}, \texttt{\detokenize{s1.p4}}, \texttt{\detokenize{s2.p1}}, \texttt{\detokenize{s3.p1}}; \texttt{\detokenize{gpt-4.1-mini}}: \texttt{\detokenize{s1.p1}}.
  \emph{Criterion:} \texttt{\detokenize{gpt-4.1-mini}}.\texttt{\detokenize{trajectory_smape}} $-$ 0.6007 $\leq 0.05$.
  \emph{Prediction:} $[-0.05, 0.05]$ (論文 Table 1 の 1 回実行の平均値。前回の再現での差は +0.001 で、137 件平均の標準誤差は約 0.024 だった).
  \emph{Observed:} pending.
  \emph{Verdict:} pending.
\end{enumerate}
\noindent\emph{Assumed, so that the claims together imply H1:}
\begin{itemize}
\item 同名モデルの現行版（Gemini-2.5-Pro / Flash の GA 版、GPT-4.1 / 4.1-mini の 2025-04-14 スナップショット）が、論文で使われた版と同等の能力を持つこと（c1〜c4）
\item aarch64 向けの libroadrunner 2.7.0 と antimony 2.14.0 が、論文環境の libroadrunner 2.8.0 と同じ軌道と評価値を与えること（c1〜c4）
\item temperature 1.0 での 1 回実行の平均 STE（137 件平均）の実行間変動が 0.05 より十分小さいこと（c1〜c4）
\item Vercel AI Gateway の OpenAI 互換エンドポイント経由の各モデルが、各社の公式 API と同じ応答分布を持つこと（c1〜c4）
\item 平均 STE が Table 1 の再現性を代表すること。RMS（modifier あり / なし）と NTS は同じ実行から表として報告するが、判定基準には用いない（c1〜c4）
\end{itemize}
\noindent\textbf{Sources.}
\noindent\textbf{S1} \texttt{\detokenize{duan-2025-measuring}}: Measuring Scientific Capabilities of Language Models with a Systems Biology Dry Lab (2025).
\begin{itemize}
\item[\texttt{\detokenize{s1.p1}}] (setup) ``We evaluated six models from three frontier model families, each represented by their professional and mini variants: Anthropic’s Claude-3.7-Sonnet-20250219 (P) and Claude-3.5-Haiku-20241022 (M); Google’s Gemini-2.5-Pro-Preview-03-25 (P) and Gemini-2.5-Flash-Preview-04-17 (M); and OpenAI’s GPT-4.1-2025-04-14 (P) and GPT-4.1-Mini-2025-04-14 (M).''
\item[\texttt{\detokenize{s1.p2}}] (setup) ``Due to computational resource constraints, we restricted our evaluation to SCIGYM-small, with a total of 137 tasks.''
\item[\texttt{\detokenize{s1.p3}}] (setup) ``The agent may submit its final model and conclude the discovery process at any point, subject to a maximum iteration limit (20 in our benchmark).''
\item[\texttt{\detokenize{s1.p4}}] (setup) ``If the submitted model is invalid or cannot be simulated, we allow for 3 additional debugging iterations before evaluating against the incomplete SBML structure.''
\item[\texttt{\detokenize{s1.p5}}] (result) ``Gemini-2.5-Flash 0.4181 0.1527 0.1071 0.1217 0.2399 0.1839 0.2005 GPT-4.1-mini 0.6007 0.1516 0.1253 0.1320 0.2530 0.2313 0.2322 Claude-3.5-Haiku 0.6281 0.0858 0.0421 0.0530 0.1454 0.0805 0.0987 Gemini-2.5-Pro 0.3212 0.2138 0.1664 0.1817 0.3781 0.3219 0.3383 GPT-4.1 0.4611 0.2067 0.1597 0.1740 0.3517 0.2888 0.3038 Claude-3.7-Sonnet 0.3615 0.1780 0.1698 0.1688 0.3160 0.3170 0.3047''
\item[\texttt{\detokenize{s1.p6}}] (claim) ``Our results show that across all model families, pro variants consistently outperform their mini counterparts, with both lower STE errors and higher RMS scores.''
\item[\texttt{\detokenize{s1.p7}}] (claim) ``Among all models evaluated, Gemini-Pro achieved the best overall performance.''
\item[\texttt{\detokenize{s1.p8}}] (definition) ``We consider two reactions "matched" if they have identical sets of reactants and products. In the second, more stringent version, we additionally require matching modifiers.''
\item[\texttt{\detokenize{s1.p9}}] (definition) ``We employ the Symmetric Mean Absolute Percentage Error (SMAPE) (Makridakis, 1993), which normalizes errors relative to magnitude.''
\item[\texttt{\detokenize{s1.p10}}] (result) ``Table 1: Pro models outperform their mini counterparts in SCIGYM, with Gemini dominating in both size categories. We report the average Simulation Trajectory Error (STE) and Reaction Matching Score (RMS) across all benchmark instances.''
\item[\texttt{\detokenize{s1.p11}}] (definition) ``Network Topology Score (NTS). To assess success in recovering the correct system topology, we measure pairwise interactions between all species in the system and calculate F1 score. If identical relationships between the same species appear in multiple reactions, we count them only once.''
\end{itemize}
\noindent\textbf{S2} \texttt{\detokenize{h4duan-2025-scigym}}: h4duan/SciGym (2025).
\begin{itemize}
\item[\texttt{\detokenize{s2.p1}}] (setup) ``temperature: 1.0''
\item[\texttt{\detokenize{s2.p2}}] (setup) ``max\_iterations: 5''
\item[\texttt{\detokenize{s2.p3}}] (setup) ``eval\_debug\_rounds: 5''
\end{itemize}
\noindent\textbf{S3} \texttt{\detokenize{airasorg-2026-airas}}: airas-org/airas-eval (2026).
\begin{itemize}
\item[\texttt{\detokenize{s3.p1}}] (method) ``SCIGYM\_COMMIT = "88a7b93609e35b6ecb4eb343d816d6ff09256c6a"''
\end{itemize}
